\section{Every negative-integer second jet of the bilateral quadratic spectrum}
\label{hqj:sec:second}

For $0<a<1$, $b>0$, define
\begin{equation}\label{hqj:eq:E}
E(s;a,b)=\sum_{n\in\Z}\bigl((n+a)^2+b^2\bigr)^{-s},\qquad \Rea s>\tfrac12.
\end{equation}
Its Poisson representation is
\begin{align}
E(s;a,b)=\frac1{\Gamma(s)}\biggl\{&\sqrt\pi\Gamma(s-\tfrac12)b^{1-2s}
\notag\\
&+4\pi^s b^{1/2-s}\sum_{n\ge1}n^{s-1/2}\cos(2\pi na)
K_{s-1/2}(2\pi bn)\biggr\}.
\label{hqj:eq:Poisson}
\end{align}
For completeness, insert
$Q^{-s}=\Gamma(s)^{-1}\int_0^\infty t^{s-1}e^{-Qt}\dd t$
in \eqref{hqj:eq:E}. Gaussian Poisson summation converts the lattice
heat kernel into
$\sqrt{\pi/t}\sum_{k\in\Z}e^{-\pi^2k^2/t}e^{2\pi ika}$.
The zero mode integrates to the first term in braces. Pair the positive
and negative modes, and use the standard integral
$\int_0^\infty t^{\nu-1}e^{-\beta t-\gamma/t}\dd t
=2(\gamma/\beta)^{\nu/2}K_\nu(2\sqrt{\beta\gamma})$
to obtain the second. Initial absolute convergence justifies these
steps, and the exponentially convergent Bessel series then continues
the result \cite{bessel}.

The expression in braces is holomorphic at every $s=-m$, $m\ge0$.
Since $1/\Gamma(s)$ has a simple zero there,
\begin{equation}\label{hqj:eq:negative-zero}
E(-m;a,b)=0.
\end{equation}
Write
\[
L_m(a,b)=E'(-m;a,b),\qquad M_m(a,b)=E''(-m;a,b).
\]
The first jet $L_m$ is the finite family in the predecessor. The second
jet differentiates the Bessel order and cannot be obtained merely by
squaring or differentiating its already specialized value.

\subsection{An explicit half-integer order derivative}
For $m\ge0$ let
\begin{equation}\label{hqj:eq:PAB}
P_m(y)=\sum_{j=0}^m\frac{(m+j)!}{j!(m-j)!2^j}y^j,
\qquad A_m(y)=(-1)^mP_m(-y).
\end{equation}
Define polynomials $B_m$ by
\begin{equation}\label{hqj:eq:Brec}
B_0(y)=0,\quad B_1(y)=2y,\quad
B_{m+1}=B_{m-1}+(2m+1)yB_m+2yP_m\quad(m\ge1).
\end{equation}
All have degree at most $m$. Let
$E_1(x)=\int_x^\infty e^{-t}\dd t/t$ for $x>0$.

\begin{lemma}[Finite polynomial form of the Bessel order derivative]
\label{hqj:lem:Bessel}
For $x>0$,
\begin{equation}\label{hqj:eq:Korder}
\left.\partial_\nu K_\nu(x)\right|_{\nu=m+1/2}
=\sqrt{\frac\pi{2x}}e^{-x}
\left\{A_m(1/x)e^{2x}E_1(2x)+B_m(1/x)\right\}.
\end{equation}
At $\nu=-m-1/2$ the derivative has the opposite sign.
\end{lemma}
\begin{proof}
The classical base is $K_{1/2}(x)=k(x)=\sqrt{\pi/(2x)}e^{-x}$.
Set $f(x)=e^{2x}E_1(2x)$, so
$f'=2f-1/x$ and $f''-2f'=1/x^2$.
The modified Bessel equation, differentiated in its order at $\nu=1/2$,
is
\[
x^2D''+xD'-(x^2+\tfrac14)D=k.
\]
Substituting $D=kf$ reduces it exactly to
$x^2(f''-2f')=1$. The real integral representation of $K_\nu$ and
Laplace localization at its minimum give
$(\partial_\nu K_\nu)/k=O(x^{-1})$ for fixed $\nu$ as $x\to\infty$;
this also follows from the differentiated standard large-$x$ expansion.
Our $f$ has the same decay. Their difference solves the homogeneous
Bessel equation: the growing solution is excluded by decay, and the
remaining multiple of $k$ is excluded by its ratio tending to zero.
Thus \eqref{hqj:eq:Korder} holds at $m=0$.

Differentiate the order recurrence
$K_{\nu+1}=K_{\nu-1}+2\nu K_\nu/x$.
The evenness $K_{-\nu}=K_\nu$ implies oddness of its order derivative.
At $\nu=1/2$, this gives the $m=1$ polynomial $B_1=2y$ and
$A_1=y-1$. Inductively the differentiated recurrence gives
\eqref{hqj:eq:Brec}. The coefficient of $e^{2x}E_1(2x)$ satisfies the
same homogeneous recurrence and equals $(-1)^mP_m(-y)$ by the finite
recurrence for $P_m$. This proves the formula at all $m$.
\end{proof}
The Bessel recurrence and order differentiability are classical
\cite{besselrec,order}; this lemma records a convenient exact
normalization for the spectral calculation. The first polynomials are
\begin{center}
\begin{tabular}{c l l l}
\toprule
$m$ & $P_m(y)$ & $A_m(y)$ & $B_m(y)$\\
\midrule
0 & $1$ & $1$ & $0$\\
1 & $1+y$ & $y-1$ & $2y$\\
2 & $1+3y+3y^2$ & $1-3y+3y^2$ & $2y+8y^2$\\
3 & $1+6y+15y^2+15y^3$ & $-1+6y-15y^2+15y^3$
& $4y+16y^2+46y^3$\\
\bottomrule
\end{tabular}
\end{center}

\subsection{One explicit polylogarithm integral family}
For $p\in\C$ set
\begin{equation}\label{hqj:eq:J}
\mathcal J_p(a,b)=\int_1^\infty
\frac{\Li_p(e^{-2\pi bw+2\pi ia})}{w+1}\dd w.
\end{equation}
This is an ordinary absolutely convergent integral, entire in $p$.
For real $a$ and $b>0$, its series form is
\begin{equation}\label{hqj:eq:Jseries}
\mathcal J_p(a,b)=\sum_{n\ge1}
\frac{e^{2\pi ian+2\pi bn}E_1(4\pi bn)}{n^p}.
\end{equation}
Indeed, expand the polylogarithm under the integral and substitute
$t=2\pi bn(w+1)$ in each exponential integral. For $p$ in a compact
set, each $p$ derivative only adds a power of $\log n$; the factor
$e^{-2\pi b n w}$ dominates this uniformly when $b$ is bounded below
by a positive constant. This proves normal convergence and all claimed
interchanges, including parameter differentiation on smaller strips in
complex $a$.

Two exact identities make the added coordinate compatible with the
existing calculus:
\begin{equation}\label{hqj:eq:J-calculus}
\partial_a\mathcal J_p=2\pi i\mathcal J_{p-1},\qquad
\partial_b\mathcal J_p=2\pi\mathcal J_{p-1}
-\frac{\Li_p(e^{-2\pi b+2\pi ia})}{b}.
\end{equation}
The first follows termwise. For the second, differentiate
\eqref{hqj:eq:Jseries} and use
$\frac{\dd}{\dd b}E_1(4\pi bn)=-e^{-4\pi bn}/b$.
Every anchored primitive in $a$ is therefore explicit; for example
\[
\int_{a_0}^{a_1}\mathcal J_p(a,b)\dd a
=\frac{\mathcal J_{p+1}(a_1,b)-\mathcal J_{p+1}(a_0,b)}{2\pi i}.
\]
The integral family is introduced to give a complete expression. No
claim that it cannot be eliminated by another identity is made.

\subsection{The complete second-jet formula}
Let $q=e^{-2\pi b+2\pi ia}$ and, for real $p$, abbreviate
\[
U_p=\Rea\Li_p(q),\qquad V_p=\Rea\Li_p^{[1]}(q),\qquad
I_p=\Rea\mathcal J_p(a,b).
\]
Write $P_m(y)=\sum_jp_{mj}y^j$, $A_m(y)=\sum_ja_{mj}y^j$,
$B_m(y)=\sum_jb_{mj}y^j$. Define
\begin{equation}\label{hqj:eq:Lzero}
L_m^{(0)}(b)=-\frac{2^{2m+1}(m!)^2\pi}{(2m+1)!}b^{2m+1},
\qquad c_m(b)=\gamma-H_m+\log(\pi/b).
\end{equation}

\begin{theorem}[All nonpositive-integer second jets]\label{hqj:thm:second}
For every $m\ge0$, $0<a<1$, and $b>0$,
\begin{align}
M_m(a,b)={}&4L_m^{(0)}(b)
\bigl(H_{2m+1}-H_m-\log(2b)\bigr)\notag\\
&+\frac{4(-1)^m m!b^m}{\pi^m}
\sum_{j=0}^m\frac1{(2\pi b)^j}
\bigl\{p_{mj}(c_m(b)U_{m+j+1}-V_{m+j+1})\notag\\[-2pt]
&\hspace{43mm}-b_{mj}U_{m+j+1}-a_{mj}I_{m+j+1}\bigr\}.
\label{hqj:eq:second-all}
\end{align}
Thus a finite set of polylogarithms, their first order derivatives, and
$\mathcal J_1,\ldots,\mathcal J_{2m+1}$ suffices for each second jet.
\end{theorem}
\begin{proof}
Write \eqref{hqj:eq:Poisson} as $E(s)=H(s)/\Gamma(s)$.
The reciprocal Gamma expansion at $s=-m+\varepsilon$ is
\[
\frac1{\Gamma(-m+\varepsilon)}
=(-1)^mm!\left\{\varepsilon+(\gamma-H_m)\varepsilon^2
+O(\varepsilon^3)\right\}.
\]
Hence
\begin{equation}\label{hqj:eq:second-H}
M_m=2(-1)^mm!\bigl(H'(-m)+(\gamma-H_m)H(-m)\bigr).
\end{equation}
For the zero mode, its logarithmic derivative is
$\psi(-m-1/2)-2\log b$. The recurrence for digamma gives
\[
\psi(-m-\tfrac12)=-\gamma-2\log2+2H_{2m+1}-H_m.
\]
Substitution in \eqref{hqj:eq:second-H} yields the first line of
\eqref{hqj:eq:second-all}.

For a nonzero Fourier mode, differentiating
$\pi^s b^{1/2-s}n^{s-1/2}$ gives $\log(\pi n/b)$.
At $s=-m$, $K_{s-1/2}=K_{m+1/2}$ has the finite polynomial $P_m$.
Its order derivative has the minus sign in Lemma~\ref{hqj:lem:Bessel}.
The resulting mode coefficient, apart from its common prefactor, is
\[
e^{-2\pi bn}n^{-m-1}
\left\{P_m((2\pi bn)^{-1})(c_m(b)+\log n)
-B_m((2\pi bn)^{-1})
-A_m((2\pi bn)^{-1})e^{4\pi bn}E_1(4\pi bn)\right\}.
\]
Now
$\sum q^n\log n/n^p=-\Li_p^{[1]}(q)$, and the last part sums by
\eqref{hqj:eq:Jseries}. Expanding the three finite polynomials gives
\eqref{hqj:eq:second-all}. The exponentially decaying modes and their
order derivatives converge normally near $s=-m$, justifying the
preceding differentiations and summations.
\end{proof}

The lowest case is already a nontrivial exact identity:
\begin{align}
M_0(a,b)={}&8\pi b\bigl(\log(2b)-1\bigr)\notag\\
&+4\left\{(\gamma+\log(\pi/b))\Rea\Li_1(q)
-\Rea\Li_1^{[1]}(q)-\Rea\mathcal J_1(a,b)\right\}.
\label{hqj:eq:M0}
\end{align}
The formula includes the entire constant Fourier mode. It is not just
an oscillatory correction to an unevaluated spectral quantity.

An independent Hurwitz representation is available from the convergent
small-mass expansion
\begin{equation}\label{hqj:eq:Ebinomial}
E(s;a,b)=\sum_{j\ge0}\frac{(-1)^j(s)_j}{j!}b^{2j}
\bigl\{\zeta(2s+2j,a)+\zeta(2s+2j,1-a)\bigr\},
\quad b<\min(a,1-a).
\end{equation}
Taylor differentiation at $s=-m$ therefore gives an identity between
\eqref{hqj:eq:second-all} and an ordinary convergent series of Hurwitz
jets. In the code this route is implemented independently of the Bessel
polynomials and is used to test $m=0,1,2,3$.
